\section{Experiments} 
\label{sec:experiment}
\subsection{Datasets}
%\label{sec:data}
\subsubsection{Stanford Sentiment Treebank (SST2)}
The Stanford Sentiment Treebank~\cite{socher2013recursive} is a common benchmark classification dataset that contains movie reviews, to be classified according to their sentiment. On an average the reviews are about 200 words long. The dataset contains train, dev, test splits and 5 sentiment labels (very positive, positive, neutral, negative, very negative) 
\subsubsection{Books Intent Classification Dataset}
We use a proprietary dataset of around 60k annotated utterances of users interacting with a popular digital assistant about books related functionality. Each utterance is manually labeled with the user intent e.g., search for a book, read a book and others, with the task being to predict the user intent from a set of 20 intents. We choose 50k utterances as our training data, 5k as test data and 5k as dev data.
\subsection{Neural Models used for Compression}
We evaluate our compression strategy on two types of neural models that are commonly used for classification tasks:
%To test our compression strategy we test 2 different architectures, the first is the DAN architecture~\cite{iyyer2015deep} and the second is a LSTM based classfication architecture~\cite{joulin2016bag}. 
\subsubsection{LSTM}
Long-Short term memory neural networks (LSTMs)~\cite{hochreiter_lstm}, are powerful and popular models for classification tasks. The LSTM unit is an RNN extension, that is designed to handle long term dependencies through the use of an input gate, an output gate, a forget gate and a memory cell. Given a sentence input $\overline{X}=x_1, \cdots x_T $ and corresponding word embedding input vectors $e_i, i=1, \cdots T$, the LSTM computes a representation $r_t$ at each $t$, which is denoted as $r_t = \phi ( e_t ,r_{t-1} )$. Here, for sentence classification, we collect the full sentence representation obtained from the final timestep of the LSTM and pass it through a softmax for classification:
\begin{gather}
r^{sent} = r^{forw}_{T}, \space  
\hat{S}=softmax( W_{s} r^{sent} + b_{s})
\end{gather}
where $r^{sent}$ is the sentence representation, and $\hat{S}$ is the label predicted for the sentence. For our classification experiments, we use 300-dimensional Glove ~\cite{pennington2014glove} embeddings as input. Our vocabulary size $V$ contains only the words that appear in the training data, and we compress the input word embedding matrix $W$ (original size $V \times 300$).
\subsubsection{DAN}
The Deep Averaging Network (DAN) model proposed ~\cite{iyyer2015deep} is a bag-of-words neural model that averages the word embeddings in each input utterance as the sentence representation $r^{sent}$. $r^{sent}$ is passed through a series of fully connected layers, and fed into a softmax layer for classification. Assume an input sentence  $\overline{X}$ of length $T$, and corresponding word embeddings $e_i$, then the sentence representation is:  
\begin{gather}
r^{sent} = \frac{1}{T} \sum_{i=1}^T { e_i }
\end{gather}
DAN has been proven a state-of-the-art model for simple classification tasks such as the Factoid Question Answering task~\cite{iyyer2015deep}. In out setup, we have two fully connected layers after the sentence representation $r^{sent}$, of sizes 1024 and 512 respectively. For the DAN experiments we tried DAN-RAND \cite{iyyer2015deep} variant where we take a randomly initialized 300 dimensional embedding. For this experiment our Baseline2 is to use appropriate low rank randomly initialized embedding matrix and train on that from the beginning. Setting up the experiments with randomly initialized as well as GloVe initialized ensures that all the gains we report is not due to low dimensional latent structure of the initialization embedding matrix but are attributed to our proposed approach. 
